\begin{theorem}[Symmetric-function valuation bounds]\label{mod:ramification-symmetricbounds}
Let $K/F$ be a finite Galois extension of nonarchimedean local fields and
write
\[
 P_i(x)=\operatorname{Tr}_{K/F}(x^i),\qquad
 s_j(x)=s_j(\{\sigma(x):\sigma\in\operatorname{Gal}(K/F)\}).
\]
For an integer $D$ and a positive integer $p=[K:F]$, isolate the exact
trace-lattice input
\begin{equation}\label{eq:lean-trace-lower-bound}
 v_K(z)\ge q
 \quad\Longrightarrow\quad
 v_F(\operatorname{Tr}_{K/F}z)
 \ge \left\lfloor\frac{q+D}{p}\right\rfloor
 \qquad(q\in\mathbf Z).
\end{equation}
Then, if $v_K(x)\ge a$ and $i\ge1$,
\begin{equation}\label{eq:lean-power-sum-bound}
 v_F(P_i(x))\ge\left\lfloor\frac{ia+D}{p}\right\rfloor.
\end{equation}
This is the public theorem \texttt{powerSum\_bound}.  The trace-lattice
equality itself belongs to the different and trace-ideal nodes; the present
node takes precisely its lower-bound half
\eqref{eq:lean-trace-lower-bound} as an explicit hypothesis, rather than
assuming an undeclared theorem.

For an arbitrary finite family $z=(z_\sigma)$ in a commutative ring, Lean
first proves the algebraic Newton identity
\[
 j s_j(z)=(-1)^{j+1}
 \sum_{\substack{r+i=j\\r<j}}(-1)^r s_r(z)P_i(z).
\]
After evaluating this identity on the Galois conjugates and using the
existing \texttt{Extension.elementarySymmetric} API, Lean retains this
antidiagonal form; equivalently, after reindexing, it is
\begin{equation}\label{eq:lean-newton-symmetric}
 j s_j(x)=\sum_{i=1}^{j}(-1)^{i-1}s_{j-i}(x)P_i(x),
 \qquad s_0(x)=1.
\end{equation}
The algebraic identities and the valuation arguments are separate
declarations.

In the wild prime-degree case assume
\[
 p=\operatorname{char}k_F=[K:F],\qquad
 T\ge2,\qquad D=(p-1)T,
\]
and assume \eqref{eq:lean-trace-lower-bound} with this value of $D$.
For every $1\le j<p$, Lean proves the exact manuscript bound
\begin{equation}\label{eq:lean-wild-symmetric-bound}
 v_F(s_j(x))\ge
 \left\lfloor\frac{ja+(p-1)T}{p}\right\rfloor.
\end{equation}
The proof uses that $j<p=\operatorname{char}k_F$ makes $j$ a unit, including
an explicit residue-characteristic divisibility argument, and that
$(p-1)T\ge p$.  There is also a source-level specialization in which
$K/F$ is a prime cyclic extension, $t>0$ is an actual lower break,
$T=t+1$, and the ramification index is $p$.  It proves
\eqref{eq:lean-wild-symmetric-bound} for the entire nonterminal range and
separately proves the endpoint
\begin{equation}\label{eq:lean-symmetric-endpoint}
 v_F(s_p(x))=v_F(N_{K/F}x)=v_K(x)\ge a.
\end{equation}

For the tame odd prime-degree range, let $\ell=[K:F]$ be prime,
$\ell\ne2$, $\operatorname{char}k_F\ne\ell$, and
$e(K/F)=\ell$.  If $2\le j<\ell$, then the direct conjugate-monomial
argument gives
\begin{equation}\label{eq:lean-tame-symmetric-bound}
 v_F(s_j(x))\ge
 \left\lceil\frac{ja}{\ell}\right\rceil
 =\left\lfloor\frac{ja+\ell-1}{\ell}\right\rfloor.
\end{equation}
The proof explicitly transports the estimate upstairs using
$v_K|_F=\ell v_F$; it does not use the wild Newton estimate.

The norm-truncation specialization retains exactly the range
$2\le j<p$: for $q,r\ge0$, if
\[
 pr\le2q+(p-1)T,
\]
then $v_K(x)\ge q$ implies $v_F(s_j(x))\ge r$ for every
$2\le j<p$.  Thus the range is empty when $p=2$, while $j=1$ and $j=p$
remain the trace and norm terms.

Finally put
\[
 e_0=\left\lfloor\frac{(p-1)T}{p}\right\rfloor
     =T-\left\lceil\frac{T}{p}\right\rceil.
\]
The trace of $1$ gives
$e_0\le v_F(p)$, and for odd $p$ one has $T\le2e_0$.  The exact numerical
inequalities used for the negative-valuation reciprocal terms are also
proved: for $c\ge0$,
\[
 e_0\le(p-2)c+
 \left\lfloor\frac{(p-1)c+(p-1)T}{p}\right\rfloor,
\]
and, for $2\le j<p$,
\[
 e_0\le(j-2)c+
 \left\lfloor\frac{(p-j)c+(p-1)T}{p}\right\rfloor.
\]
These retain every floor and every endpoint needed in the later proof of
Proposition~\ref{prop:X-depth}.  The completed finite-products API remains
the source of the prime-field identities used by those later correction
calculations; none is duplicated here.
\LeanFile{LanglandsFirstMainLemma/Ramification/SymmetricBounds.lean}
\lean{LanglandsFirstMainLemma.FiniteFamily.finitePowerSum}
\lean{LanglandsFirstMainLemma.FiniteFamily.finiteElementarySymmetric}
\lean{LanglandsFirstMainLemma.FiniteFamily.newton}
\lean{LanglandsFirstMainLemma.galoisPowerSum}
\lean{LanglandsFirstMainLemma.algebraMap_galoisPowerSum_eq_sum_conjugates}
\lean{LanglandsFirstMainLemma.elementarySymmetric_newton_identity}
\lean{LanglandsFirstMainLemma.TraceIdealLowerBound}
\lean{LanglandsFirstMainLemma.integerCeilingDiv}
\lean{LanglandsFirstMainLemma.powerSum_bound}
\lean{LanglandsFirstMainLemma.degree_natCast_ord_bound}
\lean{LanglandsFirstMainLemma.ord_natCast_eq_zero_of_lt_residueCharacteristic}
\lean{LanglandsFirstMainLemma.elementarySymmetric_degree_ord}
\lean{LanglandsFirstMainLemma.elementarySymmetric_degree_bound}
\lean{LanglandsFirstMainLemma.totallyRamified_elementarySymmetric_bound}
\lean{LanglandsFirstMainLemma.tame_elementarySymmetric_bound}
\lean{LanglandsFirstMainLemma.wildDifferentContribution}
\lean{LanglandsFirstMainLemma.wild_elementarySymmetric_bound}
\lean{LanglandsFirstMainLemma.wild_symmetric_bound}
\lean{LanglandsFirstMainLemma.wild_elementarySymmetric_uniform_bound}
\lean{LanglandsFirstMainLemma.wild_intermediateSymmetric_bound}
\lean{LanglandsFirstMainLemma.wildBaseDepth}
\lean{LanglandsFirstMainLemma.integerCeilingDiv_eq_neg_ediv}
\lean{LanglandsFirstMainLemma.wildBaseDepth_eq_sub_integerCeilingDiv}
\lean{LanglandsFirstMainLemma.two_mul_wildBaseDepth_ge}
\lean{LanglandsFirstMainLemma.wild_reciprocal_first_depth}
\lean{LanglandsFirstMainLemma.wild_reciprocal_intermediate_depth}
\leanok
\SourceText{Lemma~\ref{lem:wild-symmetric-bound} and Proposition~\ref{prop:X-depth}.}
\uses{mod:ramification-primecyclicextension,mod:localfield-extension,mod:localfield-valuation,mod:basic-finiteproducts}
\FormalizationContract The different contribution is exactly
$D=(p-1)(t+1)$, the trace-lattice lower bound is an explicit input, and
the wild proof uses the exact floor estimate rather than an asymptotic
substitute.  The tame proof uses the exact ceiling and the ramification-index
scaling upstairs.  The ranges $1\le j<p$, $2\le j<p$, and the endpoint
$j=p$ are represented by separate declarations; every theorem that uses a
denominator assumes its positivity or primality, and the residue-characteristic
condition making $j$ a unit is explicit.
\end{theorem}
